% Vecteur accélération d'une voiture en trois points d'un virage circulaire (exercice 5, chapitre 13)
\begin{tikzpicture}[x=1cm,y=1cm]
\draw[otRule,line width=7mm] (195:2.6) arc[start angle=195,end angle=355,radius=2.6];
\draw[otMuted,line width=0.4pt,dash pattern=on 3pt off 2pt] (195:2.6) arc[start angle=195,end angle=355,radius=2.6];
\fill[otInk] (0,0) circle (0.05) node[pcLegende,anchor=south] {centre du virage};
\draw[otMuted,line width=0.5pt,-{Stealth[length=1.8mm]}] (300:3.25) arc[start angle=300,end angle=330,radius=3.25];
\node[pcLegende,text=otMuted,anchor=west] at (315:3.35) {sens du parcours};
\foreach \ang/\nom/\dev in {215/A/0,270/B/-28,325/C/28}{%
  \fill[otInk] (\ang:2.6) circle (0.07);
  \node[pcLegendeG,anchor=center] at (\ang:3.1) {\nom};
  \draw[pcAcc] (\ang:2.6) -- ++({\ang+180+\dev}:1.25);}
\end{tikzpicture}
\par\smallskip\legende{Le vecteur accélération de la voiture en A, B et C (même échelle).}
